% ============================================================================
\section{GPAS: Spending Computation on Noisy Task Estimates}
\label{sec:theory}
% ============================================================================

\subsection{Measure fluctuations where they affect the update}

GPAS asks how much a task's micro-batch gradient varies around its own mean.
A large gradient shared by every micro-batch is a consistent learning signal.
A large fluctuation between micro-batches is something additional samples
can average out. We measure the latter after the optimizer's coordinate
scaling, so the statistic reflects the coordinates in which the update
will move.

Let $D$ be the diagonal scaling available at the start of the step.
For SGD, $D=I$; for Adam-type optimizers,
$D=\operatorname{diag}((\sqrt{\widehat v_{\rm pre}}+\epsilon)^{-1})$,
where $\widehat v_{\rm pre}$ is the pre-step bias-corrected second moment.
Define the task's noise and the combined-gradient variance as
\begin{equation}
    e_i=\mathbb E\|D(g_i-\mu_i)\|_2^2,
    \qquad
    V(m)=\mathbb E\|D(A-\mu)\|_2^2
         =\sum_{i=1}^M\frac{w_i^2e_i}{m_i},
    \label{eq:batch_variance}
\end{equation}
where $\mu=\sum_iw_i\mu_i$. The equality follows from conditionally
independent micro-batches at a fixed training state. It captures a simple
benefit of averaging: doubling a task's count halves its contribution to
the variance. We use the pre-step $D$ as a local measure for AdamW and leave
the optimizer update unchanged.

The minimum-variance continuous allocation is
\begin{equation}
    m_i\propto w_i\sqrt{e_i}.
    \label{eq:batch_gpas_allocation}
\end{equation}
This is the classical Neyman allocation
\citep{neyman1934representative,cochran1977sampling}, applied to MOPD
micro-batches in the optimizer's coordinates. For equally weighted tasks,
nine times the noise calls for three times as many micro-batches. The square
root reflects diminishing returns: as a task receives more samples, its
mean becomes more precise and further samples help less. We scale the counts
to sum to $G$, enforce the bounds, and round by largest remainder.
Appendix~\ref{app:allocation_proofs} gives the derivation and
Appendix~\ref{app:descent} connects the variance to a local descent bound.

\subsection{Reuse the gradients already computed}
\label{sec:online_estimates}

After processing a task's micro-batches, let $\bar g_i$ be their sample mean.
The observed noise is
\begin{equation}
    \widehat e_i=\frac1{m_i-1}\sum_{s=1}^{m_i}
                    \|D(g_{i,s}-\bar g_i)\|_2^2.
    \label{eq:noise_estimate}
\end{equation}
We smooth these measurements across steps,
$\widetilde e_i\leftarrow\rho\widetilde e_i+(1-\rho)\widehat e_i$,
and use $\widetilde e_i$ in the next allocation. Smoothing lets individual
long or unusual responses influence the schedule gradually. The first step
uses Uniform counts and initializes the moving averages.

Tasks are processed in blocks. A running mean and a scalar sum of squared
deviations give Equation~\ref{eq:noise_estimate} without retaining every
micro-batch gradient. The same parameter-sized mean buffer is reused for
each task. The additional work consists of vector reductions and the count
calculation; rollout, teacher scoring, and backward passes are the ordinary
training passes. We measure the memory and running-time cost of collecting
these statistics.

This measurement also makes the schedule easy to interpret. We log task
loss, response length, measured noise, and assigned counts together.
Loss tracks the remaining distillation gap; noise tracks how variable
the estimate of its update is. GPAS can therefore give two tasks with
similar losses different counts, or change their relative counts as
the optimizer's scaling evolves. These trajectories connect the allocation
to the training process directly, without requiring a separate token-level
noise model.

\subsection{Account for the time a micro-batch takes}

Tasks also differ in cost. A long code response may occupy the rollout
engine and teacher service much longer than a short instruction-following
response. GPAS has a cost-aware mode that uses measured time alongside
gradient noise.

Let $\tau_i$ be the marginal time for a task-$i$ micro-batch and $C$ the
fixed time for synchronization, the optimizer, and other step-level work.
For a step-time model $T(m)=C+\sum_i m_i\tau_i$, the scheduler minimizes
\begin{equation}
    J_C(m)=T(m)V(m)
          =\left(C+\sum_i m_i\tau_i\right)
           \sum_i\frac{w_i^2e_i}{m_i}.
    \label{eq:batch_cost_objective}
\end{equation}
This local criterion balances precision per update against updates per
unit time. With four tasks, we enumerate the feasible integer count vectors.
When fixed time is negligible, the continuous rule becomes
$m_i\propto w_i\sqrt{e_i/\tau_i}$: an expensive task needs more variability
to justify the same number of micro-batches. As fixed step time grows,
the choice approaches the per-step rule. We fit the timing model to the
training trace and evaluate efficiency using end-to-end GPU hours.

\begin{algorithm}[t]
\caption{One training step with GPAS.}
\label{alg:batch_gpas}
\small
\begin{minipage}{0.97\linewidth}
\begin{enumerate}
    \item Choose counts $m$ from the previous noise and time averages;
    use Uniform on the first step. Read the pre-step optimizer scaling $D$.
    \item For each task, generate $m_i$ micro-batches and obtain its
    teacher's top-$k$ scores on the student responses.
    \item Compute the ordinary dense loss. Accumulate the task's
    gradients, record their variation under $D$, and add the task mean
    to $A$ with weight $w_i$.
    \item Update the noise and time averages, then apply the usual
    optimizer step with gradient $A$.
\end{enumerate}
\end{minipage}
\end{algorithm}
